\documentclass{article}
\usepackage{writeright}

\title{A Well-Formed Sample Paper}
\author{Alice Smith and Bob Jones}

\begin{document}

\begin{abstract}
This paper presents a study of academic writing policies. We propose a framework for
automated enforcement of writing conventions in LaTeX documents.
\end{abstract}

\section{Introduction}
\label{sec:intro}
We investigate the problem of academic writing policy enforcement in Section~\ref{sec:method}.
Our method automatically detects violations at compile time.
Natural language processing (NLP) tools assist in language checks.
We evaluate our approach on a corpus of academic papers.

\section{Method}
\label{sec:method}
We describe our approach in this section; see also Section~\ref{sec:results}.
The framework consists of two components: a Chrome extension and a LaTeX package.

\begin{figure}[t]
  \centering
  \includegraphics[width=0.8\linewidth]{example-image}
  \caption{Architecture of the WriteRight system.}
  \label{fig:architecture}
\end{figure}

As shown in Figure~\ref{fig:architecture}, the system has two main parts.

\begin{table}[t]
  \centering
  \caption{Comparison of policy enforcement methods.}
  \label{tab:comparison}
  \begin{tabular}{lcc}
    \hline
    Method & Speed & Accuracy \\
    \hline
    Manual & Slow & High \\
    Automated & Fast & Medium \\
    \hline
  \end{tabular}
\end{table}

Table~\ref{tab:comparison} shows performance comparisons.

\section{Results}
\label{sec:results}
We evaluate our method on 50 papers; see Section~\ref{sec:conclusion} for discussion.

\begin{equation}
  F_1 = \frac{2 \cdot P \cdot R}{P + R}
  \label{eq:f1}
\end{equation}

We compute the $F_1$ score using Equation~\ref{eq:f1}.

\section{Conclusion}
\label{sec:conclusion}
We presented WriteRight, a system for enforcing academic writing policies.
This extends the work of Section~\ref{sec:intro}.

\begin{thebibliography}{9}
\bibitem{knuth}
  Donald Knuth, \textit{The \TeX book}, Addison-Wesley, 1984.
\end{thebibliography}

\end{document}
